% ---------------------------------------------------------------------------
% Limitations and Responsible AI (part of Section 6) - drafted for M3.
% Follows error_analysis.tex, which opens Section 6 with the Phase 8 analysis.
% Citation keys match report/references.bib.
% ---------------------------------------------------------------------------
\subsection{Limitations}
\label{sec:limitations}
\begin{enumerate}
  \item \emph{Unknown speakers.} The dataset gives no speaker identities, so the frozen split cannot be
        speaker-disjoint. The same person may appear in both training and test clips, and every score may overstate
        performance on new speakers. Keyword-spotting benchmarks usually avoid this by keeping each speaker in one
        partition~\cite{warden2018speech}.
  \item \emph{A closed vocabulary.} Every clip is forced into one of 12 words, and there is no ``unknown word'' or
        ``silence'' class. A deployed system would confidently mislabel anything outside the vocabulary.
  \item \emph{One region, one collection.} All recordings come from crowd-workers in Kenya, collected once. The
        results may not transfer to other varieties of Swahili (for example in Tanzania or the DRC), to other
        recording devices, or to conversational speech.
  \item \emph{Augmentation constrained by the rules.} External data such as noise corpora is not allowed. Our
        synthetic noise (10--30~dB SNR) was harsher than most real clips (median SNR 42~dB), so the augmentation
        study is only a partial test.
  \item \emph{Selection on one validation split.} Configurations were chosen from many validation runs, so
        validation scores are optimistic~\cite{cawley2010overfitting}. Differences of about one seed standard
        deviation (0.004 log loss for A3, 0.023 for A4) are indistinguishable from noise. The test split, never used
        for selection, gives the unbiased estimate.
  \item \emph{Early stopping.} A patience of five epochs stopped one seed each of A3 and A4 at epoch 16 of 40, while
        the learning rate was still high. For A4 that run explains most of the larger seed spread; A3's seeds stayed
        close (0.004). We did not retune the patience after seeing the seed runs, to avoid selecting on validation
        once more.
  \item \emph{Label and collection noise.} 10 test and 11 validation clips defeat every order-aware model. Confident
        learning and an external recogniser identify one likely label error, three words said in English, and two
        recordings where the crop missed the word. For most of the rest the recogniser hears other speech or
        nothing, which is weak evidence on its own (Section~\ref{sec:errors}). No one listened to the clips: our
        evidence is automatic, and the recogniser is itself fallible on short, noisy clips. We did not relabel anything, so these clips count as errors for every
        model. Removing the clear cases changes test accuracy by only 0.3 points.
  \item \emph{Pretraining confounds.} The two A5 backbones differ in size (316M against 95M parameters) as well as
        in pretraining language, and XLS-R's pretraining included only 91 hours of Swahili. A5's gains therefore
        cannot be attributed to Swahili data or to the language of pretraining alone. A Transformer trained from
        scratch, which would separate the architecture from the pretraining, was not run.
  \item \emph{A small test set.} With 630 test clips, accuracy is known to within about $\pm$1--3 points, and
        each sound-alike pair has only about 105 clips.
  \item \emph{Sensitive metrics.} A few confident errors can dominate log loss, so it is always reported with
        accuracy and calibration error. ECE in turn depends on the binning, and its low-confidence bins are sparse.
\end{enumerate}

\subsection{Responsible AI}
\paragraph{Consent and data protection}
Voice recordings are biometric data. The clips were crowdsourced by the challenge organiser under its own terms,
and we honoured its rules: the audio and derived features were never published, were shared only through private
datasets among the group, and are excluded from the public repository.

\paragraph{Fairness we could not measure}
The data carries no speaker metadata (gender, age, dialect or region), so we could not test whether error rates
differ between groups. The wide pitch range suggests a mix of voices, but that is no substitute for a fairness
audit. Speakers of other Swahili varieties, children and older adults may be served worse by any model trained
only on these recordings.

What we could measure is recording conditions, and they matter. In noisy recordings (SNR below 30~dB), the
order-aware models make 8--13\% errors, against at most 1.5\% in cleaner ones (Section~\ref{sec:errors}). Users
with cheap phones or in noisy places, often the same users that spoken interfaces aim to include, would therefore
see several times more errors than the headline accuracy suggests.

\paragraph{Consequences of errors}
The intended applications, such as spoken digits for mobile money, helplines or menu navigation for users with
limited literacy~\cite{doumbouya2021using}, turn a misrecognised digit into a wrong transaction or a wrong service.
Calibration therefore matters as much as accuracy. A deployed system should confirm uncertain inputs, which is why
we report log loss and calibration error, and why we apply temperature scaling. Because the vocabulary is closed,
it also needs a way to reject speech outside the vocabulary.

\paragraph{Dual use}
Keyword spotting in under-resourced African languages has humanitarian uses, such as monitoring radio for relief
programmes~\cite{menon2018fast}. The same technology could also be used for surveillance of speech. The models
here recognise only 12 isolated words, which limits that risk.

\paragraph{Environmental cost}
The neural models were trained on one Kaggle Tesla~T4, except one A4 run on a laptop CPU; the classical models ran
on CPUs. All 19 logged A3 runs took about 10 minutes of GPU time in total. All of A4's runs, including the failed
and superseded ones, took under 30 minutes. The pretrained A5 dominates the budget: its 16 runs, including one failed
run and the three data-size runs, took about 2.8 hours. That is about 11 minutes per run, against about one minute
or less for A3 and A4. Even so, the compute and energy cost of the study is small, and we fine-tuned the smallest
XLS-R model rather than the 1B or 2B versions.

\paragraph{Use of AI assistance}
We used AI to plan the project and to help us understand some complex situations.
